\begin{abstract}
When should auxiliary data from another embodiment be used, in what proportion, and how does that decision change with data scale? We study these questions through the optimal mixing risk of standard full-ridge regression in a dual-domain spectral model with task mismatch.
Our main result is a fixed-dimension expansion of the actual full-ridge risk and its first two derivatives. For a population risk with a unique interior minimizer and positive curvature, we prove a first-order sample-size correction to the finite-sample optimal mixing weight, together with first-order expansions of the optimized risk and gain. These expansions allow us to analyze the global optimum over all mixing weights in our setting, rather than inferring global mixing decisions from the zero-mixing derivative alone. A counterexample shows why the local derivative alone cannot determine that optimum.
Under matched covariances, we derive a first-order expansion of the optimized full-ridge gain and an exchange formula based on the first-order risk approximation. An explicitly defined quadratic surrogate admits an exact piecewise dosage rule, with an interior hyperbolic branch and explicit boundary handling. Alignment, variance reduction, and spectral coverage provide the mechanism explanation. Direct $n_r$ tests pass pre-registered root-compatibility checks, and additional experiments examine small nonlinear learners and real CIFAR features. The results remain limited to controlled statistical models: general non-unimodal classification, general-spectrum or tuned-regularization extensions, the proportional regime, and robot-policy validation remain open.
\end{abstract}
